\documentclass[10pt,a4paper]{amsart}
\usepackage[margin=19mm]{geometry}
\usepackage{amsmath,amssymb,mathtools}
\usepackage[scr=rsfs,cal=euler]{mathalfa}
\usepackage{xfrac}
\usepackage[unicode,hidelinks,bookmarks=false]{hyperref}
\newcommand{\ie}{i.e.\ }
\newcommand{\cf}{cf.\ }
\newcommand{\resp}{respectively\ }
\newcommand{\Subsection}{\subsection}
\providecommand{\nobreakdash}{\nobreak\hbox{-}}
\setlength{\emergencystretch}{2em}
\multlinegap0pt
\usepackage{amssymb}
\usepackage{xcolor}
\newtheorem{lem}{Lemma}
\newtheorem{thm}{Theorem}
\DeclareMathOperator{\Aut}{Aut}
\newcommand{\C}{\mathbb{C}}
\newcommand{\La}{\mathfrak{a}}
\newcommand{\Lg}{\mathfrak{g}}
\newcommand{\ad}{{\rm ad}}
\newcommand{\al}{\alpha}
\newcommand{\la}{\lambda}
\renewcommand{\phi}{\varphi}
\newcommand{\ra}{\rightarrow}
\title{Fixed-point-free automorphisms of solvable Lie algebras}
\author{Dietrich Burde, Karel Dekimpe}
\date{Source: arXiv:2604.27916v1 (CC BY 4.0)}
\begin{document}
\maketitle

\noindent\emph{Source and licence.} Fixed-point-free automorphisms of solvable Lie algebras. Version arXiv:2604.27916v1 (CC BY 4.0), distributed by arXiv under the Creative Commons Attribution 4.0 International licence, http://creativecommons.org/licenses/by/4.0/, as recorded in the arXiv licence metadata for this exact version and shown on the abstract page https://arxiv.org/abs/2604.27916v1. Full licence notice: \texttt{licenses/arxiv-2604.27916-CC-BY-4.0.txt}.\par\smallskip

\noindent\textit{Editorial scope.} Counted unit: the article's thm 5.4 (source0.tex lines 722--726). The article's complete original proof of it is delivered with it (source0.tex lines 728--834, one \texttt{proof} environment of 107 lines). No mathematics is written, repaired or re-derived: the statement and the proof are the article's own lines, in the article's own notation, and no retained line is substituted or re-flowed.  Retained uncounted as the source-local context the proof needs: lines 649--663.  Nothing was omitted from any retained span, so the package carries no omission record.  No retained line contains a citation, so the wrapper carries no bibliography and no bibliography entry.  No retained line contains a figure environment, an image-inclusion command or drawing code, so no image is delivered and the package declares no asset dependency.  Editorial formatting only, in addition: an AMS \texttt{amsart} preamble that restates the article's own declarations and macros used by the retained lines, namely the environments \texttt{lem}, \texttt{thm} on separate counters, together with the standard packages those lines need.  The retained environments print in this document as Lemma 1, Theorem 1 in order of appearance, whereas the article prints the same items with the numbering of its own full document.  The complete native source of the article as distributed in the arXiv e-print for this exact version is bundled unchanged as \texttt{sources/199509/source0.tex}.
\begin{lem}\label{5.3}
Let $B\in GL_m(\C)$ and $n\ge 2$. Then $B$ is $n$-cyclotomic if and only if there exists an
invertible matrix $A\in GL_k(\C)$ with $m=kn$, such that $B$ is similar to
\[
M_A=  
\begin{pmatrix}
0 & 0 & 0 & \cdots & 0 & A \cr
I_k & 0 & 0 & \cdots & 0 & 0 \cr
0 & I_k & 0 & \cdots & 0 & 0 \cr
%0 & 0 & I_k & \cdots & 0 & 0 \cr
\vdots & \vdots & \vdots & \ddots & \vdots & \vdots \cr
0 & 0 & 0 & \cdots & I_k & 0 
\end{pmatrix}  
\]  
\end{lem}  

\begin{thm}\label{5.4}
Let $\Lg=\La\rtimes \langle v\rangle$ be a complex almost abelian Lie algebra with $Z(\Lg)=0$. Then
$\Lg$ admits an f.p.f.\ automorphism if and only if $\ad(v)_{\mid \La}$ is $n$-cyclotomic for some
$n\ge 2$. In this case $\Lg$ admits an f.p.f.\ automorphism of order $kn$ for all $k\ge 2$.
\end{thm}  

\begin{proof}
Since we have $Z(\Lg)=0$, the linear map $\ad(v)_{\mid \La}: \La\ra \La$ has no eigenvalue zero. Hence it is a
vector space isomorphism. Assume first that $\Lg$ admits an f.p.f.\ automorphism $\phi$. We may assume that $\phi$
is semisimple. Because of $\phi(\La)\subseteq \La$ we have $\phi(v)=\al v+a_v$ with $a_v\in \La$.
It follows that $\al$ is an eigenvalue of $\phi$, so that $\al\neq 0,1$. Let
\[
\La=V_{\la_1}\oplus V_{\la_2}\oplus \cdots \oplus V_{\la_n}
\]
be the decomposition of $\La$ into a direct sum of eigenspaces $V_{\la_i}$ for $\phi$, with distinct eigenvalues
$\la_i$ of $\phi$. Fix an eigenvalue $\la_i$. Then we have $[v,V_{\la_i}]\subseteq V_{\al\la_i}$. Indeed, with
$w\in V_{\la_i}$ we have
\begin{align*}
\phi([v,w]) & =[\phi(v),\phi(w)] \\
            & = [\al v +a_v, \la_i w] \\
            & = \al\la_i[v,w],
\end{align*}
because $[a_v,w]=0$. So we have $[v,w]\in V_{\al\la_i}$ as required. As $\ad(v)_{\mid \La}$ is a vector space
isomorphism of $\La$ we have that
\[
\dim ([v,V_{\la_i}])=\dim (V_{\la_i})\le \dim (V_{\al\la_i})
\]
Hence we obtain a series of subspaces
\[
V_{\la_i}, V_{\al\la_i}, V_{\al^2\la_i},\ldots , V_{\al^k\la_i}, \ldots
\]
of $\La$ of non-decreasing dimension. Since $\La$ is the direct sum of all eigenspaces, these
subspaces cannot be all different. So there exist $p>q$ such that $\al^q\la_i=\al^p\la_i$, or
$\al^{p-q}=1$. It follows that $\al$ is an $n$-th primitive root of unity, where $n\ge 2$ because of $\al\neq 1$.
We have
\[
\dim (V_{\la_1})\le \dim (V_{\al\la_1})\le \cdots \le \dim (V_{\al^n\la_1})=\dim (V_{\la_1}),
\]  
so that we have equality everywhere, and
\[
[v,V_{\al^k\la_1}]=V_{\al^{k+1}\la_1}
\]
for all $k\ge 0$. Let $V_1=V_{\la_1}\oplus V_{\al\la_1}\oplus \cdots \oplus V_{\al^{n-1}\la_1}$. If $V_1\neq \La$,
then there exists an eigenvalue $\la_2$ of $\phi$ such that $\la_2\neq \al^k\la_1$ for all $k$. Then
$V_2=V_{\la_2}\oplus V_{\al\la_2}\oplus \cdots \oplus V_{\al^{n-1}\la_2}$ is another subspace of $\La$ with
$V_1\cap V_2=0$, $V_1\oplus V_2\subseteq \La$ and $[v,V_{\al^k\la_2}]\subseteq V_{\al^{k+1}\la_2}$ for all $k$.
Continuing this way, we find finitely many eigenvalues $\la_1,\ldots , \la_m$, each giving rise to a sequence of
subspaces $V_{\la_i}, V_{\al\la_i},\ldots ,V_{\al^{n-1}\la_i}$ with $[v,V_{\al^k\la_i}]\subseteq V_{\al^{k+1}\la_i}$ for all $k$.
Now let
\begin{align*}
U_1 & = V_{\la_1}\oplus V_{\la_2}\oplus \cdots \oplus V_{\la_m} \\
U_2 & = V_{\al\la_1}\oplus V_{\al\la_2}\oplus \cdots \oplus V_{\al\la_m} \\
\vdots \hspace*{0.18cm} & = \hspace*{0.2cm} \vdots \\
U_n & = V_{\al^{n-1}\la_1}\oplus V_{\al^{n-1}\la_2}\oplus \cdots \oplus V_{\al^{n-1}\la_m} \\
\end{align*}
Note that all spaces $U_i$ have equal dimension, say $\dim (U_i)=k$. Then we have
\begin{align*}
U_1\oplus U_2\oplus \cdots \oplus U_n & = \La, \\
[v,U_i] & = U_{i+1},\; \text{for all} \; i<n,\\
[v,U_n] & = U_1.
\end{align*}  
Choose a basis $\{u_{11},\ldots ,u_{1k}\}$ of $U_1$, and let $u_{2i}=[v,u_{1i}]$ for $i=1,\ldots ,k$.
Then $\{u_{21},\ldots ,u_{2k}\}$ is a basis of $U_2$. Continuing this way, with $u_{j+1,i}=[v,u_{ji}]$ for $j<n$
we obtain a basis of each subspace $U_j$ for $1 \le j\le n$. With respect to the basis
\[
\{u_{11},\ldots ,u_{1k},u_{21},\ldots ,u_{2k}, u_{31},\ldots ,u_{nk}\}
\]  
of $\La$ the matrix of $\ad(v)_{\mid \La}$ is given by
\[
M_A=  
\begin{pmatrix}
0 & 0 & 0 & \cdots & 0 & A \cr
I_k & 0 & 0 & \cdots & 0 & 0 \cr
0 & I_k & 0 & \cdots & 0 & 0 \cr
\vdots & \vdots & \vdots & \ddots & \vdots & \vdots \cr
0 & 0 & 0 & \cdots & I_k & 0 
\end{pmatrix}  
\]
for some $A\in GL_k(\C)$. By Lemma $\ref{5.3}$ it follows that $\ad(v)_{\mid \La}$ is $n$-cyclotomic. \\[0.2cm]
Conversely, assume that  $\ad(v)_{\mid \La}$ is $n$-cyclotomic. Then there is a basis $\{u_{11},u_{12}, \ldots ,u_{nk}\}$
of $\La$ such that the matrix of $\ad(v)_{\mid \La}$ with respect to this basis is given by $M_A$ as above. Now
choose a $\la\in \C^{\times}$ and a primitive $n$-th root of unity $\zeta$ such that
\[
\zeta^{\ell}\la\neq 1 \text{ for all } 0\le \ell \le n-1.
\]
According to the form of $M_A$, we can decompose $\La$ into subspaces $U_i$ such that
\begin{align*}
U_1\oplus U_2\oplus \cdots \oplus U_m & = \La, \\
[v,U_i] & = U_{i+1},\; \text{for all} \; i<n,\\
[v,U_n] & = U_1.
\end{align*}
We define a map $\phi\colon \Lg\ra \Lg$ for $\Lg=\La\rtimes \langle v\rangle$ by $\phi(v)=\zeta v$, and
$\phi_{\mid U_i}$ as the multiplication with $\zeta^{i-1}\la$. So with respect to the above basis, $\phi_{\mid \La}$
is represented by the block matrix
\[
{\rm diag}(\la I_k,\zeta \la I_k,\ldots ,\zeta^{n-1}\la I_k).
\]
Since $\zeta^{\ell}\la\neq 1$, it follows that $\phi$ is fixed-point-free. We claim that
$\phi\in \Aut(\Lg)$. For this it is enough to show that for
each $1\le i\le n$ and for each $u\in U_i$
we have
\[
\phi([v,u])=[\phi(v),\phi(u)].
\]
Since $[v,u]\in U_{i+1}$, with $U_{n+1}=U_1$, we have $\phi([v,u])=\zeta^i\la [v,u]$, and
\[
[\phi(v),\phi(u)]=[\zeta v,\zeta^{i-1}\la u]=\zeta^i\la [v,u],
\]
which was to be shown. \\[0.2cm]
Finally, if $\ad(v)_{\mid \La}$ is $n$-cyclotomic for some $n\ge 2$, we can choose 
$\la=\zeta_{kn}$ in the definition of $\phi$ for every $k\ge 2$. Then
$\phi$ is an f.p.f.\ automorphism of order $kn$. 
\end{proof}  

\end{document}
